% Generated draft: requires mathematical and source review.
\documentclass[11pt]{article}
\usepackage[margin=25mm]{geometry}
\usepackage{fontspec}
\setmainfont{Latin Modern Roman}
\usepackage{amsmath,amssymb,mathtools,unicode-math}
\setmathfont{Latin Modern Math}
\usepackage{enumitem,needspace,xurl,hyperref}
\setlength{\parindent}{0pt}
\setlength{\parskip}{5pt}
\setlength{\emergencystretch}{3em}
\newcommand{\N}{\mathbb{N}}
\newcommand{\Z}{\mathbb{Z}}
\newcommand{\Q}{\mathbb{Q}}
\newcommand{\R}{\mathbb{R}}
\newcommand{\C}{\mathbb{C}}
\newcommand{\F}{\mathbb{F}}
\newcommand{\A}{\mathbb{A}}
\newcommand{\PP}{\mathbb{P}}
\newcommand{\E}{\mathbb{E}}
\newcommand{\eps}{\varepsilon}
\newcommand{\dd}{\,\mathrm{d}}
\begin{document}
\section*{UnsolvedMath ID: 800645}
\textbf{Source problem code:} catalogue entry
\subsection{Definitions and mathematical statement}
For a topological space \(X\), an open cover is a family of open subsets whose union is \(X\). Under the usual convention, an \(\omega\)-cover is an open cover \(\mathcal U\) such that \(X\notin\mathcal U\) and every finite subset of \(X\) is contained in some member of \(\mathcal U\). Write \(\mathcal O(X)\) and \(\Omega(X)\) for the corresponding families of covers, subject to the source conventions that remain unspecified below.

For families \(\mathcal A\) and \(\mathcal B\) of collections of subsets of \(X\), the selection principle \(\mathsf S_1(\mathcal A,\mathcal B)\) means that for every sequence \((\mathcal U_n)_{n\geq 1}\) with \(\mathcal U_n\in\mathcal A\), there exist choices \(U_n\in\mathcal U_n\) such that \(\{U_n:n\geq 1\}\in\mathcal B\).

Let \(\mathcal G(X)\) denote the cover family written in the supplied record as \(\mathcal O^{\omega\text{-}*\mathrm{gp}*}(X)\). This is a placeholder for the source notation: the supplied extract does not define this family or explain whether the asterisks are mathematical notation or extraction artifacts. No definition of \(\mathcal G(X)\) is imposed here.

The question in Problem 18 is whether, for every space \(X\) in the source's intended class of spaces,
\[
\mathsf S_1\bigl(\Omega(X),\Omega(X)\bigr)
\quad\Longleftrightarrow\quad
\mathsf S_1\bigl(\Omega(X),\mathcal G(X)\bigr).
\]
The intended class of spaces, any restriction to countable covers, and the precise conventions for \(\mathcal O(X)\) are not specified in the supplied extract. Consequently, this is a conditional formulation requiring those source definitions before it becomes a fully specified mathematical statement.

As background, the extract attributes to the reference labeled \([\mathrm{strongdiags}]\) the assertion that \(\mathsf S_1(\Omega,\mathcal G)\) is strictly stronger than \(\mathsf S_1(\Omega,\mathcal O)\), and states that \(\mathsf S_1(\Omega,\mathcal O)\) is equivalent to \(\mathsf S_1(\mathcal O,\mathcal O)\). These are supplied contextual assertions, retained with their unresolved ambient hypotheses.
\subsection{Short English statement}
The problem asks whether selecting one set from each of a sequence of \(\omega\)-covers to obtain an \(\omega\)-cover is equivalent to obtaining the specially designated cover family in the source. The supplied record labels this question open. Its special cover definition and ambient hypotheses are missing, so the formulation requires human review.
\subsection{Sources}
\begin{itemize}
\item Local UnsolvedMath registry, permanent ID \(800645\), source problem code \texttt{catalogue entry}, AMR Open Problem Lists. The registry attributes the record to Tsaban, gives proposed year 2003, and labels its status open; this status is supplied metadata, not an independently verified literature assessment.
\item Tsaban, \emph{Selection Principles in Mathematics: A Milestone of Open Problems} (2003), Problem 18. Supplied source URL: \url{https://arxiv.org/abs/math/0312182}. The record reports access on 2026-10-16 and extraction by \texttt{arxiv\_latex}; no independent source visit or verification is claimed here.
\item The supplied attribution/context is \([\mathrm{coc8},\text{ Problem }4]\). The extract also cites \([\mathrm{strongdiags}]\) for its strict-strength comparison. Bibliographic details for these two reference labels are not supplied.
\end{itemize}
\end{document}
